\begin{abstract}
We propose two parameter-efficient architectural additions to hybrid state-space/attention language models --- \emph{hormone routing}, a token-conditional router that injects a gated combination of frozen direction vectors into the residual stream after each attention block, and \emph{intracellular attention}, a slot-attention mechanism over the $\dstate$ memory cells of a selective state-space block --- and use them to investigate the relationship between architectural inductive bias and the optimisation objective. Both mechanisms use a zero-initialised output gate so that the unmodified hybrid baseline is recovered exactly at initialisation. We pretrain a 125M-parameter Mamba/MQA hybrid on FineWeb-Edu under a fixed 2B-token budget and report a striking decoupling between the routing module's representational organisation and its effect on the loss objective. The hormone-routing module develops semantically-coherent per-genre routing specialisation across every variant we examine --- including controls that replace the extracted directions with random unit vectors, with vectors extracted from a randomly-initialised model, or that freeze the output gate at $1.0$. Maximum pairwise $L_1$ distance between genre routing distributions ranges from $0.6$ to $1.6$ (uniform routing $= 0.0$, disjoint $= 2.0$), and the variant with the \emph{highest} specialisation strength is the one with a frozen gate that the model cannot modulate. Yet \textbf{specialisation strength does not correlate with loss}: the four properly-initialised variants converge to identical validation perplexity ($26.13$), and the variant with the second-highest specialisation strength (\textsc{HR-noext}, $L_1 \approx 0.73$) is the only one that meaningfully changes the loss --- and changes it in the wrong direction ($+2.07$ PPL). The model is invariant to the routing module's \emph{content} (extracted, random, or fixed-gate) but not invariant to its \emph{training-schedule timing} (introducing it during base pretraining versus as a continuation). We interpret this as evidence that hybrid language models can develop structured representational organisation along architectural channels that the optimisation objective does not see --- a phenomenon in tension with the implicit assumption that interpretable internal structure should track downstream task performance. Both contributions add fewer than $1\%$ to the base model's parameter count and are released as open-source code.
\end{abstract}
